% 本文件由 LaTeX 排版 Agent 根据有效 translation.json 自动生成。
% author=Ignatius of Antioch; work_id=0108; title=依纳爵致斐拉德尔菲亚人书

\chapter{\WorkTitle{依纳爵致斐拉德尔菲亚人书}}

% structure: p0001 source=h1 target=omitted reason=page-title-already-rendered-by-parent

% structure: p0002 source=h2 target=section reason=relative-heading-rank; base=h2

\section{问候}

依纳爵，又名\OriginalTerm{德奥福鲁}{Theophorus}，致在亚洲的斐拉德尔菲亚的教会，即天父天主及我等主耶稣基督的教会。此教会已蒙怜悯，在天主的和谐中建立，因我等主的苦难而不息欢悦，并藉祂的复活充满一切怜悯。我以耶稣基督的血问候你们，祂是我们永恒不竭的喜乐，尤其当你们与主教、长老及执事合而为一时。他们乃是按照耶稣基督的旨意被选派，由祂按自己的旨意并藉圣神圣定在稳固之中。\par

% structure: p0004 source=h2 target=section reason=relative-heading-rank; base=h2

\section{第一章 称赞主教}

这位主教，我知道，他获得那关乎公共利益的职分，并非出于自己，也不是藉着人\BibleRef{迦 1:1}，也不是因虚荣，而是由于天父天主及主耶稣基督的爱。他的温良令我惊异，他沉默所能成就的，胜过那些空谈之人。因为他与天主的诫命和谐，如同竖琴与琴弦。因此，我的灵魂称他的心意在天主前为有福的，知道它是良善而完美的，并且他的稳固与毫无怒气，乃是效法永生天主无限温良的榜样。\par

% structure: p0006 source=h2 target=section reason=relative-heading-rank; base=h2

\section{第二章 维持与主教的合一}

所以，作为光明和真理的子女，要逃避分裂和邪恶的教义；但牧人在哪里，你们就跟随如羊。因为有许多貌似可信的豺狼，藉着有害的快乐，掳掠\BibleRef{弟后 3:6}那些奔向天主的人；但在你们的合一中，它们将无立足之地。\par

% structure: p0008 source=h2 target=section reason=relative-heading-rank; base=h2

\section{第三章 避免分裂者}

要远离那些邪恶的植物，它们是耶稣基督所不栽培的，因为不是父的栽植。并非我发现了你们中间有分裂，而是极其纯洁。因为凡属于天主和耶稣基督的，也都与主教同在。凡藉着补赎回头进入教会合一的，这些人也必属于天主，好使他们按照耶稣基督而生活。我的弟兄们，不要被迷惑。若有人跟随那在教会中制造分裂的人，他必不能承受天主的国。若有人随从异样的意见，他就不认同基督的苦难。\par

% structure: p0010 source=h2 target=section reason=relative-heading-rank; base=h2

\section{第四章 只应有一个圣体圣事等}

所以要当心，只应有一个圣体圣事。因为只有一个我等主耶稣基督的身体，和一只杯，以彰显祂血的合一；一个祭台；如同只有一位主教，连同长老团和执事，我的同仆。好使你们无论做什么，都按天主的旨意而行。\par

% structure: p0012 source=h2 target=section reason=relative-heading-rank; base=h2

\section{第五章 请为我祈祷}

我的弟兄们，我在爱你们中大大扩张；极其喜乐，我谋求确保你们的安全。然而这不是我，而是耶稣基督，我因祂被捆绑而更加惧怕，因我尚未完全。但你们的祈祷将使我完全，好使我达到那藉怜悯所分给我的份，我逃向福音如同耶稣的肉，逃向宗徒如同教会的长老团。我们也当爱先知们，因为他们也宣讲了福音，并寄希望于祂，等候祂；他们因信靠祂，藉与耶稣基督的合一而得救，是圣洁的人，值得爱戴和钦佩，由耶稣基督为他们作证，并在共同希望的福音中与我们一起被记念。\par

% structure: p0014 source=h2 target=section reason=relative-heading-rank; base=h2

\section{第六章 不要接受犹太教}

若有人向你们宣讲犹太法律，不要听从他。因为从一个受过割礼的人那里听从基督的教义，比从一个未受割礼的人那里听从犹太教更好。但若这类人中有人不谈论耶稣基督，在我眼中他们不过是死人的纪念碑和坟墓，上面只写着人的名字。所以，要逃避这世界之王的邪恶诡计和网罗，免得一旦被他的伎俩征服，你们的爱心变得软弱。但要以不可分割的心彼此结合。我感谢我的天主，因我对你们怀有良好的良心，无人能私下或公开夸口我曾在大事小事上加重任何人的负担。我愿所有我与之交谈过的人，不以此为反对他们的证据。\par

% structure: p0016 source=h2 target=section reason=relative-heading-rank; base=h2

\section{第七章 我已劝你们合一}

因为虽然有人曾按肉体欺骗我，但圣神既然是出于天主，就不受欺骗。因为祂知道从何而来，往何而去\BibleRef{若 3:8}，并察验人心的隐秘。因为当我在你们中间时，我曾呼喊，大声说：\Quote{要留意主教、长老团和执事}。现在，有些人怀疑我这样说，是预先知道你们中间某人所造成的分裂。但祂是为我而受捆绑的那位为我作证，我并未从任何人那里得到消息。但圣神宣告了这些话：\Quote{不要在没有主教的情况下做事；要保守你们的身体如同天主的殿；要爱合一；逃避分裂；要跟随耶稣基督，如同祂跟随祂的父一样}。\par

% structure: p0018 source=h2 target=section reason=relative-heading-rank; base=h2

\section{第八章 同上（续）}

因此，我做了分内之事，作为一个献身于合一的人。因为哪里有分裂和忿怒，天主就不住在哪里。主赐予所有悔改的人宽恕，若他们悔改归向天主的合一，并与主教共融。我信赖耶稣基督的恩宠，祂必使你们摆脱一切束缚。我劝你们不要出于纷争做任何事，而要按基督的教导而行。当我听到有些人说：\Quote{若我在古代圣经中找不到，我就不相信福音}；我对他们说：\Quote{经上记着}，他们回答说：\Quote{那还有待证实}。但对我而言，耶稣基督取代了一切古代的东西：祂的十字架、死和复活，以及由祂而来的信德，是无瑕疵的古代纪念碑；我藉此，因你们的祈祷，希望得以成义。\par

% structure: p0020 source=h2 target=section reason=relative-heading-rank; base=h2

\section{第九章 旧约是好的：新约更美}

司铎固然是好的，但大司铎更好；祂受托管理至圣之所，唯独祂获信托付天主的奥秘。祂是圣父的门，\Person{亚巴郎}、\Person{依撒格}、\Person{雅各伯}、众先知、宗徒及教会，皆由此门进入。这一切皆以达至天主的合一为鹄的。但福音拥有超越先前时期的特质，即我主\Person{耶稣基督}的显现、祂的苦难与复活。因为蒙爱的先知曾预言祂，但福音乃不朽的完美。这一切若你们以爱德相信，便都是好的。\par

% structure: p0022 source=h2 target=section reason=relative-heading-rank; base=h2

\section{第十章 祝贺\Place{安提约基雅}人教难结束}

既然，按照你们的祈祷和你们在\Person{基督耶稣}内所怀的怜悯，有人向我报告，在\Place{叙利亚}的\Place{安提约基雅}的教会得到了平安，那么你们作为天主的教会，理应选立一位执事，作为天主的使者［为你们］前往［那里的弟兄们］那里，好使他在他们聚会时能与他们一同喜乐，并光荣［天主的］圣名。那堪当如此职事的人，在\Person{耶稣基督}内是有福的；你们也必得光荣。如果你们愿意，为天主的缘故，此事并非超出你们的能力；就如邻近的教会已经派遣，有的派了主教，有的派了司铎和执事。\par

% structure: p0024 source=h2 target=section reason=relative-heading-rank; base=h2

\section{第十一章 感谢与问候}

至于\Place{基里基雅}的执事\Person{斐洛}，一位有声誉的人，他仍在天主的话语上服事我，连同\Person{雷乌斯·阿加托普斯}，一位被拣选的人，他从\Place{叙利亚}跟随我，不顾惜自己的性命——这些人在你们那里作见证；我也为你们感谢天主，因为你们接待了他们，如同主接待了你们。但那些羞辱他们的人，愿他们因\Person{耶稣基督}的恩宠获得宽恕！\Place{特洛阿}弟兄们的爱德问候你们；我也藉着\Person{布尔布斯}写信给你们，他是\Place{厄弗所}人和\Place{斯米纳}人派来与我同行的，以示他们的敬意。愿主\Person{耶稣基督}光荣他们，他们盼望于祂，在肉身、灵魂、信德、爱德与和好之中！在\Person{基督耶稣}内，我们的共同希望中，平安！\par

\SourceRecord{Translated by Alexander Roberts and James Donaldson.；Ante-Nicene Fathers；Vol. 1.；Edited by Alexander Roberts, James Donaldson, and A. Cleveland Coxe.；Buffalo, NY: Christian Literature Publishing Co.,；1885.；<http://www.newadvent.org/fathers/0108.htm>.}
